%% kapa: exam=1 round=30 year=2019 period=1 form=A subject=civil_law subject_ko=민법 seq=39 number=39
%% source: Q-Net 감정평가사 기출문제 (cst00309) artlSeq=5205205 file=extracted/1008/5205205/2019년 제30회 감정평가사 1차시험 문제지/2019년 제30회 감정평가사 1차시험 문제지/2019년도 제30회 감정평가사 1차 1교시 A형.hwp
%% conversion: kordoc 4.24.51 -> tools/gen1.py (auto); review=proofread-sonnet; flags=none
\begin{question}{39}
저당권에 관한 설명으로 옳은 것을 모두 고른 것은? (다툼이 있으면 판례에 따름)
\begin{framedtext}
ㄱ. 저당권이 설정된 건물이 증축된 경우에 기존 건물에 대한 저당권은 법률에 특별한 규정 또는 설정행위에서 다른 약정이 없다면, 증축되어 부합된 건물부분에 대해서도 그 효력이 미친다. \par
ㄴ. 저당부동산의 교환가치를 하락시키는 행위가 있더라도 저당권자는 저당권에 기한 방해배제청구권을 행사할 수 없다. \par
ㄷ. 저당물의 제3취득자는 그 부동산의 개량을 위한 유익비를 지출하여 가치의 증가가 현존하더라도, 그 비용을 저당물의 매각대금에서 우선적으로 상환받을 수 없다. \par
ㄹ. 채권자 아닌 타인의 명의로 저당권이 설정되었다면, 피담보채권의 실질적 귀속주체가 누구인지를 불문하고 그 효력이 인정되지 않는다.
\end{framedtext}
\begin{choices}
  \item ㄱ
  \item ㄷ
  \item ㄱ, ㄷ
  \item ㄴ, ㄹ
  \item ㄱ, ㄷ, ㄹ
\end{choices}
\end{question}
\answer{1}
